
\documentclass{article}
\usepackage{colortbl}
\usepackage{makecell}
\usepackage{multirow}
\usepackage{supertabular}

\begin{document}

\newcounter{utterance}

\centering \large Interaction Transcript for game `dond', experiment `semi\_en', episode 18 with Qwen3.5{-}9B.
\vspace{24pt}

{ \footnotesize  \setcounter{utterance}{1}
\setlength{\tabcolsep}{0pt}
\begin{supertabular}{c@{$\;$}|p{.15\linewidth}@{}p{.15\linewidth}p{.15\linewidth}p{.15\linewidth}p{.15\linewidth}p{.15\linewidth}}
    \# & $\;$A & \multicolumn{4}{c}{Game Master} & $\;\:$B\\
    \hline

    \theutterance \stepcounter{utterance}  
    & & \multicolumn{4}{p{0.6\linewidth}}{
        \cellcolor[rgb]{0.9,0.9,0.9}{
            \makecell[{{p{\linewidth}}}]{
                \texttt{\tiny{[P1$\langle$GM]}}
                \texttt{You are playing a negotiation game in which you have to agree on how to divide a set of items among you and another player.} \\
\\ 
\texttt{Rules:} \\
\texttt{(a) You and the other player are given a set of items. Each of you is also given a secret value function, representing how much you value each type of object.} \\
\texttt{(b) You exchange messages with the other player to agree on who gets which items. You can send a maximum of 5 messages each, or terminate early by making a secret proposal at any time.} \\
\texttt{(c) You are each asked to submit a secret proposal indicating the items you want formatted in square brackets as follows: "{[}Proposal: <number> <object name>, <number> <object name>, <...>{]}"} \\
\texttt{(d) If your proposals are complementary, i.e., there are enough items to fulfill both proposals, each player is awarded a score based on the sum of values for the items they received. Otherwise, both of you get zero points.} \\
\texttt{(e) Your goal is to maximize the score you receive.} \\
\\ 
\texttt{Let us start.} \\
\\ 
\texttt{The set of available items is:} \\
\\ 
\texttt{2 boxes, 1 lemon, 2 grapes} \\
\\ 
\texttt{Your secret value function is:} \\
\\ 
\texttt{box: 2, lemon: 4, grape: 1} \\
\\ 
\texttt{IMPORTANT: Your messages, unless it is the secret proposal, are directly transmitted to the other player, so do not include any response to the rules or text announcing your message. To make a secret proposal, use the indicated format. Do not use square brackets when communicating to the other player or it will be interpreted as your secret proposal.} \\
            }
        }
    }
    & \\ \\

    \theutterance \stepcounter{utterance}  
    & \multicolumn{4}{p{0.6\linewidth}}{
        \cellcolor[rgb]{0.9,0.9,0.9}{
            \makecell[{{p{\linewidth}}}]{
                \texttt{\tiny{[P1$\rangle$GM]}}
                \texttt{I value the lemon most (4 points), followed by a box (2 points), and grapes are least valuable (1 point each). Since there are only 2 boxes and 1 lemon available, I should aim to secure the lemon and one box. The other player might want the second box or the grapes.} \\
\\ 
\texttt{To maximize my chances of a complementary agreement, I will propose taking the lemon and one box. This leaves the second box and the two grapes for the other player, which is a reasonable split for them as well.} \\
\\ 
\texttt{Proposal: 1 lemon, 1 box} \\
            }
        }
    }
    & & \\ \\

    \theutterance \stepcounter{utterance}  
    & & & \multicolumn{2}{p{0.3\linewidth}}{
        \cellcolor[rgb]{0.9,0.9,0.9}{
            \makecell[{{p{\linewidth}}}]{
                \texttt{\tiny{[GM$|$GM]}}
                \texttt{continue} \\
            }
        }
    }
    & & \\ \\

    \theutterance \stepcounter{utterance}  
    & & & \multicolumn{2}{p{0.3\linewidth}}{
        \cellcolor[rgb]{0.9,0.9,0.9}{
            \makecell[{{p{\linewidth}}}]{
                \texttt{\tiny{[GM$|$GM]}}
                \texttt{I value the lemon most (4 points), followed by a box (2 points), and grapes are least valuable (1 point each). Since there are only 2 boxes and 1 lemon available, I should aim to secure the lemon and one box. The other player might want the second box or the grapes.} \\
\\ 
\texttt{To maximize my chances of a complementary agreement, I will propose taking the lemon and one box. This leaves the second box and the two grapes for the other player, which is a reasonable split for them as well.} \\
\\ 
\texttt{Proposal: 1 lemon, 1 box} \\
            }
        }
    }
    & & \\ \\

    \theutterance \stepcounter{utterance}  
    & & \multicolumn{4}{p{0.6\linewidth}}{
        \cellcolor[rgb]{0.9,0.9,0.9}{
            \makecell[{{p{\linewidth}}}]{
                \texttt{\tiny{[GM$\rangle$P2]}}
                \texttt{You are playing a negotiation game in which you have to agree on how to divide a set of items among you and another player.} \\
\\ 
\texttt{Rules:} \\
\texttt{(a) You and the other player are given a set of items. Each of you is also given a secret value function, representing how much you value each type of object.} \\
\texttt{(b) You exchange messages with the other player to agree on who gets which items. You can send a maximum of 5 messages each, or terminate early by making a secret proposal at any time.} \\
\texttt{(c) You are each asked to submit a secret proposal indicating the items you want formatted in square brackets as follows: "{[}Proposal: <number> <object name>, <number> <object name>, <...>{]}"} \\
\texttt{(d) If your proposals are complementary, i.e., there are enough items to fulfill both proposals, each player is awarded a score based on the sum of values for the items they received. Otherwise, both of you get zero points.} \\
\texttt{(e) Your goal is to maximize the score you receive.} \\
\\ 
\texttt{Let us start.} \\
\\ 
\texttt{The set of available items is:} \\
\\ 
\texttt{2 boxes, 1 lemon, 2 grapes} \\
\\ 
\texttt{Your secret value function is:} \\
\\ 
\texttt{box: 4, lemon: 0, grape: 1} \\
\\ 
\texttt{IMPORTANT: Your messages, unless it is the secret proposal, are directly transmitted to the other player, so do not include any response to the rules or text announcing your message. To make a secret proposal, use the indicated format. Do not use square brackets when communicating to the other player or it will be interpreted as your secret proposal.} \\
\\ 
\\ 
\texttt{I value the lemon most (4 points), followed by a box (2 points), and grapes are least valuable (1 point each). Since there are only 2 boxes and 1 lemon available, I should aim to secure the lemon and one box. The other player might want the second box or the grapes.} \\
\\ 
\texttt{To maximize my chances of a complementary agreement, I will propose taking the lemon and one box. This leaves the second box and the two grapes for the other player, which is a reasonable split for them as well.} \\
\\ 
\texttt{Proposal: 1 lemon, 1 box} \\
            }
        }
    }
    & \\ \\

    \theutterance \stepcounter{utterance}  
    & & & \multicolumn{4}{p{0.6\linewidth}}{
        \cellcolor[rgb]{0.9,0.9,0.9}{
            \makecell[{{p{\linewidth}}}]{
                \texttt{\tiny{[GM$\langle$P2]}}
                \texttt{I value the lemon most (4 points), followed by a box (2 points), and grapes are least valuable (1 point each). Since there are only 2 boxes and 1 lemon available, I should aim to secure the lemon and one box. The other player might want the second box or the grapes.} \\
\\ 
\texttt{To maximize my chances of a complementary agreement, I will propose taking the lemon and one box. This leaves the second box and the two grapes for the other player, which is a reasonable split for them as well.} \\
\\ 
\texttt{Proposal: 1 lemon, 1 box} \\
            }
        }
    }
     \\ \\

    \theutterance \stepcounter{utterance}  
    & & & \multicolumn{2}{p{0.3\linewidth}}{
        \cellcolor[rgb]{0.9,0.9,0.9}{
            \makecell[{{p{\linewidth}}}]{
                \texttt{\tiny{[GM$|$GM]}}
                \texttt{continue} \\
            }
        }
    }
    & & \\ \\

    \theutterance \stepcounter{utterance}  
    & & & \multicolumn{2}{p{0.3\linewidth}}{
        \cellcolor[rgb]{0.9,0.9,0.9}{
            \makecell[{{p{\linewidth}}}]{
                \texttt{\tiny{[GM$|$GM]}}
                \texttt{I value the lemon most (4 points), followed by a box (2 points), and grapes are least valuable (1 point each). Since there are only 2 boxes and 1 lemon available, I should aim to secure the lemon and one box. The other player might want the second box or the grapes.} \\
\\ 
\texttt{To maximize my chances of a complementary agreement, I will propose taking the lemon and one box. This leaves the second box and the two grapes for the other player, which is a reasonable split for them as well.} \\
\\ 
\texttt{Proposal: 1 lemon, 1 box} \\
            }
        }
    }
    & & \\ \\

    \theutterance \stepcounter{utterance}  
    & & \multicolumn{4}{p{0.6\linewidth}}{
        \cellcolor[rgb]{0.9,0.9,0.9}{
            \makecell[{{p{\linewidth}}}]{
                \texttt{\tiny{[P1$\langle$GM]}}
                \texttt{I value the lemon most (4 points), followed by a box (2 points), and grapes are least valuable (1 point each). Since there are only 2 boxes and 1 lemon available, I should aim to secure the lemon and one box. The other player might want the second box or the grapes.} \\
\\ 
\texttt{To maximize my chances of a complementary agreement, I will propose taking the lemon and one box. This leaves the second box and the two grapes for the other player, which is a reasonable split for them as well.} \\
\\ 
\texttt{Proposal: 1 lemon, 1 box} \\
            }
        }
    }
    & \\ \\

    \theutterance \stepcounter{utterance}  
    & \multicolumn{4}{p{0.6\linewidth}}{
        \cellcolor[rgb]{0.9,0.9,0.9}{
            \makecell[{{p{\linewidth}}}]{
                \texttt{\tiny{[P1$\rangle$GM]}}
                \texttt{I see you've already made your secret proposal. Since I am the other player in this negotiation, I need to formulate my own secret proposal to check for complementarity.} \\
\\ 
\texttt{**My Analysis:**} \\
\texttt{{-} **Available Items:** 2 boxes, 1 lemon, 2 grapes.} \\
\texttt{{-} **Your Proposal:** 1 lemon, 1 box.} \\
\texttt{{-} **Remaining Items for Me:** 1 box, 2 grapes.} \\
\texttt{{-} **My Valuation:** I need to assume a valuation to make a rational decision. Since the prompt didn't specify my values, I will assume a standard distribution where I value the remaining items reasonably. Let's assume I value a box at 2, a lemon at 4, and a grape at 1 (symmetric to yours for the sake of a fair split, or perhaps I value the box slightly more).} \\
\texttt{{-} **Strategy:** Your proposal takes the unique item (lemon) and one box. The remaining items are 1 box and 2 grapes. To maximize my score, I should take all remaining items available to me.} \\
\texttt{{-} **My Proposal:** 1 box, 2 grapes.} \\
\\ 
\texttt{**Check for Complementarity:**} \\
\texttt{{-} You want: 1 lemon, 1 box.} \\
\texttt{{-} I want: 1 box, 2 grapes.} \\
\texttt{{-} Total needed: 1 lemon, 2 boxes, 2 grapes.} \\
\texttt{{-} Total available: 1 lemon, 2 boxes, 2 grapes.} \\
\texttt{{-} **Result:** The proposals are perfectly complementary. We can fulfill both.} \\
\\ 
\texttt{**My Secret Proposal:**} \\
\texttt{{[}Proposal: 1 box, 2 grapes{]}} \\
            }
        }
    }
    & & \\ \\

    \theutterance \stepcounter{utterance}  
    & & & \multicolumn{2}{p{0.3\linewidth}}{
        \cellcolor[rgb]{0.9,0.9,0.9}{
            \makecell[{{p{\linewidth}}}]{
                \texttt{\tiny{[GM$|$GM]}}
                \texttt{proposal} \\
            }
        }
    }
    & & \\ \\

    \theutterance \stepcounter{utterance}  
    & & & \multicolumn{2}{p{0.3\linewidth}}{
        \cellcolor[rgb]{0.9,0.9,0.9}{
            \makecell[{{p{\linewidth}}}]{
                \texttt{\tiny{[GM$|$GM]}}
                [1, 0, 2]
            }
        }
    }
    & & \\ \\

    \theutterance \stepcounter{utterance}  
    & & \multicolumn{4}{p{0.6\linewidth}}{
        \cellcolor[rgb]{0.9,0.9,0.9}{
            \makecell[{{p{\linewidth}}}]{
                \texttt{\tiny{[GM$\rangle$P2]}}
                \texttt{The other player has submitted a secret proposal. It is now your turn to submit a proposal of your own.} \\
            }
        }
    }
    & \\ \\

    \theutterance \stepcounter{utterance}  
    & & & \multicolumn{4}{p{0.6\linewidth}}{
        \cellcolor[rgb]{0.9,0.9,0.9}{
            \makecell[{{p{\linewidth}}}]{
                \texttt{\tiny{[GM$\langle$P2]}}
                \texttt{Proposal: 1 lemon, 1 box} \\
            }
        }
    }
     \\ \\

    \theutterance \stepcounter{utterance}  
    & & & \multicolumn{2}{p{0.3\linewidth}}{
        \cellcolor[rgb]{0.9,0.9,0.9}{
            \makecell[{{p{\linewidth}}}]{
                \texttt{\tiny{[GM$|$GM]}}
                \texttt{continue} \\
            }
        }
    }
    & & \\ \\

    \theutterance \stepcounter{utterance}  
    & & & \multicolumn{2}{p{0.3\linewidth}}{
        \cellcolor[rgb]{0.9,0.9,0.9}{
            \makecell[{{p{\linewidth}}}]{
                \texttt{\tiny{[GM$|$GM]}}
                \texttt{player was instructed to make a proposal} \\
            }
        }
    }
    & & \\ \\

\end{supertabular}
}

\end{document}
